\chapter*{Executive summary}
\addcontentsline{toc}{chapter}{Executive summary}
\label{sec:exec}

\textit{What this report is.} We examined every step of \banzai{}, the LCO pipeline for the FLOYDS spectrographs,
that matters for reverberation mapping (RM). We tested the archive products (\banzai{} version 1.0.1) and the current
release (1.1.0) on a Fairall~9 campaign at Siding Spring (COJ) and an NGC\,4151 campaign at Haleakal\=a (OGG). For each step we show what the pipeline
does, where and why it fails, what we changed, and what the change achieves. All corrections are target-independent
(no knowledge of the science target is used) and are built into \banzai{} (branch \texttt{rm-improvements} of
\url{https://github.com/Alymantara/banzai-floyds}); a whole data set is reduced from the raw frames with one command
(Appendix~\ref{sec:app_install}).

\textit{Status.} \textbf{Improved v3 is our final version for now.} It reduces all 833 raw frames of the two
campaigns without failures, and we verified it step by step at both sites (Chapter~\ref{sec:sites}). Table~\ref{tab:exec_numbers}
gives the headline numbers and Table~\ref{tab:exec_issues} lists every problem we found, what v3 does about it and where
in the report it is explained. We then inspected the worst
v3 reductions and found problems in two steps we had not examined: the cosmic-ray detector and the sky model under the
object (issues 27--31, Chapter~\ref{sec:cosmics}). The fixes are in the code and verified on a sample, but the
products have not yet been re-reduced with them. Every product now carries header keywords that record the quality of
its reduction, and a flag bitmask (Section~\ref{sec:qc}). The reductions compared are defined in
Table~\ref{tab:versions_names}, and the terms used throughout (orders, visits, pulls, fringe measures) in the
Conventions (p.~\pageref{sec:conventions}).

\begin{table}[!htbp]
\centering
\caption{Headline numbers. Fairall~9 at COJ and NGC\,4151 at OGG, 6\,arcsec slit. Scatters are robust rms (see
Conventions) over all exposures. ``1.1.0 as released'' is \banzai{}~1.1.0 run from raw without the object wavelength
zero point of this work (elsewhere in the report the 1.1.0 run includes it; Table~\ref{tab:versions_names}).}
\label{tab:exec_numbers}
\small
\begin{tabular}{@{}>{\raggedright\arraybackslash}p{5.6cm}cccc@{}}
\toprule
Quantity (unit) & \banzai{}~1.0.1 & \banzai{}~1.1.0 & \multicolumn{2}{c}{improved v3 (final)} \\
 & (archive) & as released & COJ & OGG \\
\midrule
{[O\,{\sc iii}]} wavelength scatter, red / blue order (\AA) & 6.4 / 2.8 & 6.6 / 3.1 & 0.69 / 0.45 & 0.94 / 0.53 \\
{[O\,{\sc iii}]} wavelength, red $-$ blue order (\AA) & $-5.9\pm3.0$ & $+0.4\pm3.1$ & $+0.30\pm0.79$ & $+1.39\pm0.88$ \\
Red-order fringe residual, rms (per cent) & 13.7 & 8.1 & 1.9 & 2.5$^{a}$ \\
Residual in the water band 8130--8360\,\AA, rms (per cent) & 16.1 & 8.0 & 3.2 & 2.7 \\
Red/blue [O\,{\sc iii}] flux ratio scatter (per cent) & 29.2 & 8.9 & 6.9 & 6.2 \\
Uncertainties: rms of the pull, red order (1 = correct errors) & 1.7--2.2$^{b}$ & up to 1.6 & 0.94--1.06 & 0.98--1.27 \\
Absolute flux level (factor) & $\times2.2$ too bright & correct & correct & correct$^{c}$ \\
\bottomrule
\end{tabular}
\par\smallskip\raggedright\footnotesize
$^{a}$Includes NGC\,4151's own lines, which fall in the measurement windows. The difference between the two
exposures of a visit, which removes all structure common to both, is the same as at COJ (0.66--0.69 per cent).
$^{b}$Beyond 7700\,\AA. $^{c}$Also across a factor-1.8 change in telescope throughput, thanks to two added standards
(Section~\ref{sec:sites_flux}). All absolute levels are biased high by the slit losses of the 2\,arcsec standards; the
usual RM scaling of each spectrum to the [O\,{\sc iii}] flux removes this bias.
\end{table}

\textit{The issues.} In Table~\ref{tab:exec_issues}, ``1.0.1'' marks problems of the archive products and ``1.1.0''
problems that remain in the current release. Issues are grouped by pipeline step, so their numbers are not in order. Status: \textit{solved} (fixed in v3 and verified), \textit{mitigated}
(reduced or flagged, not removed), \textit{open} (not addressed in the software), \textit{operations} (needs a change in
how data are taken). The open items and recommendations for LCO are collected in Chapter~\ref{sec:recommendations}.

{\small
\setlength{\LTcapwidth}{\textwidth}
\begin{longtable}{@{}>{\raggedright\arraybackslash}p{0.3cm}>{\raggedright\arraybackslash}p{5.4cm}>{\raggedright\arraybackslash}p{5.3cm}>{\raggedright\arraybackslash}p{1.6cm}>{\raggedright\arraybackslash}p{1.9cm}@{}}
\caption{Problems found, what the final pipeline (improved v3) does about them, and where in this report.}
\label{tab:exec_issues}\\
\toprule
\# & Problem and its effect & What v3 does & Where & Status \\
\midrule
\endfirsthead
\toprule
\# & Problem and its effect & What v3 does & Where & Status \\
\midrule
\endhead
\bottomrule
\endfoot
\multicolumn{5}{@{}l}{\textit{Wavelength calibration}}\\
1 & The arc fit uses no line between Hg\,4358 and 5461\,\AA{} (blue) and none blueward of 5461 or between 5791 and
6965\,\AA{} (red): H$\beta$/[O\,{\sc iii}] and H$\alpha$ are interpolated or extrapolated. 1.0.1: red order off by
$-8$\,\AA{} at Hg\,4916; 1.1.0: up to 0.6\,\AA. & Hg\,4916 and Ar\,6416, 6677, 6753, 6767, detected in every arc,
added to the arc fit with measured strengths; error at Hg\,4916 $\approx0.1$\,\AA, orders agree to 0.3\,\AA. &
\S\ref{sec:wl_lines}, \S\ref{sec:sites_arcs} & solved \\
2 & The Zn lamp is ``on'' in the headers but no Zn line is detected; Zn would fill the blue gap. & -- &
\S\ref{sec:wl_lines} & operations \\
3 & No per-exposure zero point for the object (1.0.1, 1.1.0): in a 6\,arcsec slit the object's position across the
slit shifts its spectrum by $\approx10$\,\AA{} per arcsec, which arcs and sky lines cannot measure. [O\,{\sc iii}]
scatters by 3--7\,\AA{} between exposures, up to $\pm17$\,\AA{} for NGC\,4151. & Zero point measured from the telluric
bands in the object's own light, transferred to the blue order: scatter 0.5--0.9\,\AA. & \S\ref{sec:wl_zp},
\S\ref{sec:sites_zp} & solved \\
4 & The absolute wavelength zero point is that of the telluric template, uncertain by $\approx\pm2$--4\,\AA. & Nothing:
it is a constant offset and does not affect the agreement between exposures or between orders. & \S\ref{sec:wl_zp} & open (harmless for RM) \\
\midrule
\multicolumn{5}{@{}l}{\textit{Order and object traces}}\\
5 & The object trace is a polynomial (1.0.1, 1.1.0): at the faint red end it is off by 0.8--1.4\,px (1--2 per cent
flux loss), and much worse for targets 30--60 times fainter than Fairall~9. & Physical model (atmospheric refraction along the
slit + an instrumental shape for each site + a small smooth correction fitted to the residuals): 0.1--0.3\,px down to 60 times fainter; templates for COJ and OGG.
& \S\ref{sec:traces_object}, \S\ref{sec:traces_model}, \S\ref{sec:sites_trace} & solved \\
6 & The position of the orders on the detector in each frame is fitted on the light of the science target, which
fails for faint targets. For the targets here it agrees with the lamp flat to 0.14\,px. & Unchanged (the 1.1.0 fit).
The lamp flat of the same block would measure it at a signal-to-noise ratio of thousands. & \S\ref{sec:traces_orders} & open (faint targets) \\
7 & The shapes of the orders on the detector come only from the sky flats in a list fixed in the code; for OGG the
newest are from 2019. & Sky flats already processed by the LCO archive can also be used; the 2019 OGG set works for
2024--25 data. & \S\ref{sec:traces_orders}, \S\ref{sec:sites_data} &
mitigated \\
\midrule
\multicolumn{5}{@{}l}{\textit{Fringing and extraction}}\\
8 & 1.0.1 divides every frame by a whole eight-month-old lamp flat, including its slit illumination: flux biased by up
to $\pm10$ per cent, wavelength and frame dependent. & No 2-d division by a lamp flat. & \S\ref{sec:fringe_banzai} &
solved \\
9 & 1.0.1 shifts that flat along the slit by up to hundreds of pixels (the fit locks on to wrong
matches of the repeating fringe pattern); the invalid edge of the shifted flat destroys data in 13 of 50 frames
(spikes, striped sky). & As 8; the destroyed pixels are never created. & \S\ref{sec:fringe_banzai},
\S\ref{sec:traces_object} & solved \\
10 & A lamp flat gives the phase of the fringes but not their amplitude or shape (the slit-filling lamp smooths them
differently from a point source), so scaling a lamp pattern leaves spiky residuals. & A two-harmonic fringe model with
the lamp phase, fitted to each extracted spectrum. & \S\ref{sec:fringe_harmonic} & solved \\
11 & 1.1.0 still leaves 8 per cent rms of fringing (Fairall~9). & 1.9 per cent (EG\,21: 1.5). & \S\ref{sec:raw_compare} &
solved \\
12 & Residual fringes at 6500--7500\,\AA: the phase of a single 6\,arcsec lamp flat is noisy there and slow structure
in the lamp pattern dominates it. & Phase from a phase master (median of the lamp flats of 1--2
months) after a high-pass filter; the model starts at 6200\,\AA{} and its amplitude and phase may change every
200\,\AA. & \S\ref{sec:fringe_pipeline}, \S\ref{sec:fringe_joint} & solved \\
13 & Fringes inside the telluric bands were not fitted but interpolated across. This left a 4 per cent residual in
the transmission measured on the standards, and so in every science spectrum. A 0.7 per cent residual that is fixed on
the detector also remained. & Fringes fitted jointly with a telluric template; a template of the fixed residual is
measured on the standards for each slit and removed. &
\S\ref{sec:fringe_joint} & solved \\
\midrule
\multicolumn{5}{@{}l}{\textit{Telluric correction}}\\
14 & 1.1.0 divides almost the whole red order by the standard's observed/expected ratio, so the standard's noise and
residuals (1--3 per cent) appear in every science spectrum. The transmission of a single standard is noisy. &
Transmission applied only inside the absorption bands (between the points where the absorption is 1 per cent deep),
median of the nine nearest standards. &
\S\ref{sec:flux_tell} & solved \\
15 & The saturated O$_2$ A band leaves $\pm20$ per cent residuals in 6\,arcsec spectra, because the 2\,arcsec standards
resolve it differently. The sky lines fill the slit, so they cannot give the resolution of the object's spectrum. &
Flagged in science spectra (bit 2 of the mask, 7580--7700\,\AA); smoothing the standards' transmission to match gave
no real gain. & \S\ref{sec:flux_tell} & mitigated \\
16 & The water absorption changes from night to night ($\approx3.5$ per cent in the band level); a median of standards
does not follow it. & -- & \S\ref{sec:raw_compare}, \S\ref{sec:recommendations} & open \\
17 & Beyond 8900\,\AA{} there is no telluric-free continuum: the transmission and the sensitivity are extrapolated
there (a step of $\approx30$ per cent near 9285\,\AA{} in every version). & Water band included in the telluric bands
(8900--9900\,\AA); the red end should still be used with care. & App.~\ref{sec:app_v1}, \S\ref{sec:recommendations} & open \\
\midrule
\multicolumn{5}{@{}l}{\textit{Flux calibration}}\\
18 & 1.0.1 does not scale by the exposure time: 600\,s science calibrated with 300\,s standards is 2.2 times too bright
in the archive. & Fixed in 1.1.0; v3 agrees with 1.1.0 to 4 per cent. & \S\ref{sec:flux_banzai} & solved (1.1.0) \\
19 & Standards are chosen by block or proposal, not slit width, and a single standard transfers its clouds (grey levels
of nine consecutive standards span a factor 2) and noise to the science. The 1.0.1 sensitivity functions contain the
fringes and spikes. & Median of the sensitivities of the nine nearest standards, each first scaled by its grey factor; overall level
from the clearest half; red/blue [O\,{\sc iii}] scatter 29.2 (1.0.1), 8.9 (1.1.0) $\rightarrow$ 6.9 per cent. &
\S\ref{sec:flux_banzai}, \S\ref{sec:flux_ensemble} & solved \\
20 & Nightly standards missing from the catalogue (EG\,21 at COJ; Feige\,67 and GD\,153 at OGG) are reduced as science and
never calibrate anything. & Feige\,67 and GD\,153 added (references from ESO; Feige\,67 extended beyond 9200\,\AA). EG\,21
has only a 50\,\AA{} reference. & \S\ref{sec:flux_model}, \S\ref{sec:sites_flux} & solved for OGG; open for EG\,21 \\
21 & The OGG telescope throughput rose by a factor 1.8 between 2025 April 7 and May 2 (re-aluminisation of the
primary mirror), most in the blue (48 per cent more at 3400\,\AA{} than at 5800\,\AA), and all catalogue standards are
from after it: earlier frames were calibrated 1.8 times too faint ([O\,{\sc iii}] scatter 26 per cent) and up to
30 per cent off in shape. & With the added
standards the two periods agree to 7 per cent ([O\,{\sc iii}] scatter 9.4 per cent). An ensemble must not combine
standards across such a change. & \S\ref{sec:sites_flux}, \S\ref{sec:shape_stability} & solved here; open in general; operations \\
22 & The standards are 2\,arcsec, the RM science 6\,arcsec: absolute levels are biased high by the standards' slit
losses. & Grey level from the clearest standards (lowest bias). RM scaling to [O\,{\sc iii}] removes it. &
\S\ref{sec:flux_banzai} & open (harmless for RM) \\
\midrule
\multicolumn{5}{@{}l}{\textit{Cosmic rays and the sky model} (found after the improved v3 run; Chapter~\ref{sec:cosmics})}\\
27 & The cosmic-ray detector of 1.1.0, LA Cosmic (1.0.1 has none), takes the core of a bright star in good seeing
for cosmic rays: up to 46 per cent of the core of BD$+28^\circ4211$ at OGG, whose spectrum then scatters by 2.9 times
its errors. & LA Cosmic is given a model of the object (7-pixel median along the dispersion) besides the sky; the
flagged share of the core drops from 46 to 0.1 per cent. & \S\ref{sec:cr_stars}, \S\ref{sec:cr_model} & fixed in code; re-reduction pending \\
28 & LA Cosmic clips bright narrow lines that are resolved along the slit, and the extraction scales up what is left:
H$\alpha$ of NGC\,4151 flagged in 7 of 60 frames, 2.64 and 1.18 times too bright in two. [O\,{\sc iii}]\,5007 never
flagged. Can affect any bright emission-line object in good seeing. & Object model (flagged pixels on H$\alpha$ drop from
179 to 15 and from 119 to 0 in those two frames; fluxes within 2 per cent). Wavelength bins that lose more than half of
their profile weight to flags are masked instead of scaled up, and counted in the header. & \S\ref{sec:cr_lines}, \S\ref{sec:cr_guard} & mitigated in code;
re-reduction pending \\
29 & When the object is wide (poor seeing) and the sky bright (Moon), the spline that models the sky gets a knot at
every pixel beside the object and follows the noise. The extracted flux then scatters by 3--4 times its errors (six
Fairall~9 frames), and LA Cosmic flags the object ($\approx240$ false detections per frame). Also in 1.1.0. & A sky
line is only treated as such if it stands out by 5$\sigma$ from the noise of a first-pass sky fit; the scatter drops
from 3.0--3.3 to 1.1--1.2 times the errors. & \S\ref{sec:cr_sky} & fixed in code; re-reduction
pending \\
30 & LA Cosmic (1.1.0) misses real cosmic rays on bright cores and line peaks (66 per cent found above 10$\sigma$; none
of 8 on the core of BD$+28^\circ4211$) and spreads its detections into bright regions. & With the object model, 97 per cent of
artificial cosmic rays added to the frames are found, at the cost of 1--5 more small false detections per frame on
bright traces. After LA Cosmic, the two exposures of each visit are now compared, and a flag is removed when the
partner exposure shows the same structure, which is therefore real (10--22 per cent fewer flagged pixels; on the H$\alpha$ peak
15 to 0 pixels), with no change in the agreement of their fluxes. Standards are single exposures and use the object
model only. & \S\ref{sec:cr_verify}, \S\ref{sec:cr_pair} & fixed in code; re-reduction
pending \\
31 & A warm pixel of the OGG camera at $(x,y)=(1465,409)$ is not in the bad-pixel mask (flagged as a cosmic ray in every
frame, on the wing of [O\,{\sc iii}]\,4959 of NGC\,4151). & -- & \S\ref{sec:cr_lines} & open \\
32 & Nothing in the products records how good a reduction is. & Quality keywords (\texttt{L1Q}) and a flag bitmask
in every product
(added to all 317 products of the final reductions), and a command that tabulates them. & \S\ref{sec:qc} & solved \\
\midrule
\multicolumn{5}{@{}l}{\textit{Uncertainties and products}}\\
23 & Uncertainties: 1.0.1 underestimates the red-order errors by a factor 1.7--2.2 (the fringe residuals are not in
them), 1.1.0 by up to 1.6. & Correct per pixel at both sites (rms of the pull 0.9--1.1; 1.3 at OGG where the signal is
high). Quantities integrated over many pixels need errors $\approx20$--30 per cent larger, because the fringe and
telluric residuals are not in the error array. & \S\ref{sec:raw_errors},
\S\ref{sec:sites_verify} & solved per pixel; documented \\
24 & The combined spectrum of the \texttt{-1d} product mixes the two orders, which are not aligned, over
4540--5760\,\AA. & Use the separate spectrum of each order (\texttt{EXTRACTED} table of the \texttt{-1d} product). & \S\ref{sec:recommendations} & open \\
25 & A processed 2-d frame is 47\,MB, of which 43\,MB is an uncompressed table of all pixels: a decade of data
would not fit on disk. & After a run, the diagnostic images and the pixel table are removed from the 2-d products
($\approx2$\,GB per site and season of $\approx400$ frames). & \S\ref{sec:sites_disk} & mitigated \\
26 & Installation fails on Apple-silicon laptops with an Intel Python or Intel compiler settings. & A native arm64 Python
environment, and a wrapper script that runs the pipeline without the user's compiler settings. & App.~\ref{sec:app_install} & solved \\
\end{longtable}
}

\textit{How to read the rest.} Chapter~\ref{sec:summary} summarises the pipeline and the changes step by step
(Fig.~\ref{fig:flow}, Table~\ref{tab:stages}); Chapter~\ref{sec:intro} introduces the problem and the data;
Chapter~\ref{sec:diagnostics} defines the measurements;
Chapters~\ref{sec:traces}--\ref{sec:flux} treat traces, wavelengths, fringing and flux calibration in turn;
Chapter~\ref{sec:raw} runs the pipeline from raw and compares all versions; Chapter~\ref{sec:sites} describes the final
version and verifies it at both sites; Chapter~\ref{sec:cosmics} treats cosmic rays and the sky model and defines the
quality keywords; Chapter~\ref{sec:recommendations} gives the recommendations.
